Seis ejes (E1--E6), los del Anexo~B. No responden la misma pregunta. E1 mira si OpenAPI ya es el
formato con el que se describe REST. E2 y E3 dicen de dónde pueden salir E y S. E4 separa dominio
e interfaz. E5 y E6 quedan al borde: hablan de pruebas y de ranking, no del cruce con un SD. La
\Cref{tab:sintesis-ejes} junta el recorte. El detalle de cada paper está en el Anexo~B.

\subsection{Eje E1 --- OpenAPI como estándar}

Casas et al.\ (2021) mapean 47 artículos primarios entre 2011 y 2020. OpenAPI/Swagger queda como
el estándar de facto para describir REST de forma \textit{machine-readable}. El 68\,\% de esos
trabajos propone herramientas y el 43\,\% mejora la documentación. En ese conjunto no aparece un
modelo de negocio bancario del tipo BIAN.

Mainas et al.\ (2023) anotan contratos (\texttt{x-refersTo}, \texttt{x-kindOf},
\texttt{x-mapsTo}), pasan OpenAPI a RDF y consultan con SPARQL. El corpus supera las 300
especificaciones de SwaggerHub y no es de banca. Lo usamos como pista de representación para S,
no como definición del score.

Santos y Casas (2025), con \textit{Design Science}, catalogan 59 métricas de API Gateway
(latencia, tráfico, errores) en 68 productos. Quedan anotadas para no mezclarlas con la coherencia
documental: miden el \textit{gateway} en producción.

\subsection{Eje E2 --- Matching estructural}

El eje se apoya en un solo paper, Shvaiko y Euzenat (2005). La taxonomía separa \textit{matching}
exacto y aproximado, evidencia sintáctica, semántica o externa, y el grano elemento frente a
estructura. Es anterior a OpenAPI y a JSON Schema. Para el \StructuralScore{} tomamos de ahí el
contraste de \textit{paths} y \textit{schemas} contra \Behavior{}s y \BusinessObject{}s. La
taxonomía no trae una métrica lista para el par bancario; el cálculo se escribe en el
Capítulo~\ref{cap:metodologia}.

\subsection{Eje E3 --- Matching semántico y ontologías}

El Bayadh et al.\ (2015), Chandrasekaran y Mago (2022) y Santos Sousa et al.\ (2026) no trabajan
contratos bancarios. El primero advierte que una sola medida de alineación ontológica se queda
corta; las cuatro familias están en §\ref{sec:bases-teoricas} y no se repiten aquí. Chandrasekaran
y Mago recorren la similitud desde WordNet hasta \textit{transformers}. Santos Sousa et al.\ (2026)
listan, para \textit{embeddings} en \textit{ontology matching}, tres límites que nos importan:
supervisión, escala y correspondencias complejas. Esos tres papers son candidatos de técnica para
S y, en parte, para C.

\subsection{Eje E4 --- Microservicios e interfaces}

Schmidt y Thiry (2020) parten de 715 registros y se quedan con 27 estudios primarios (2013--2019)
sobre cómo se identifican microservicios con MDE y DDD. El dato que usamos es la separación: el
\ServiceDomain{} puede ser la referencia de negocio y la API, la capa que lo expone. En su muestra
casi no hay trabajos que integren MDE y DDD a la vez. No reportan una métrica de coherencia
contrato--BIAN. El eje solo justifica que las dos capas no se confundan.

\subsection{Eje E5 --- Testing y calidad API (borde)}

Cinco papers, y OpenAPI entra como insumo de prueba. Ehsan et al.\ (2022) cubren 16 artículos de
testing RESTful; la métrica que más se repite es \textit{code coverage}. Golmohammadi et al.\
(2024) leen 92 papers (\textit{fuzzing}, \textit{search-based}, \textit{model-based}) que generan
casos sobre endpoints ya desplegados. Otieno y Araka (2026) ponen aprendizaje automático en el
testing y señalan \textit{datasets} chicos. Tanveer et al.\ (2025) van a vulnerabilidades.
Vitorino et al.\ (2025), a datos adversariales.

Ese cuerpo es el que deja fuera del plan el testing, la seguridad y la QoS operativa
(\Cref{sec:alcances}).

\subsection{Eje E6 --- Selección y recomendación de servicios}

Mahanra Rao et al.\ (2025) revisan 25 estudios (2020--2024, PRISMA) que usan BERT o
\textit{transformers} para elegir servicios web según QoS. La representación semántica de las
descripciones se parece a lo que haríamos en S, pero el resultado es un ranking, no un contraste
contra un \ServiceDomain{}. Reportan costo de cómputo alto, \textit{datasets} cortos y poca
interpretabilidad. Queda como vecino, no como técnica del modelo.

\begin{table}[H]
  \centering
  \caption{Síntesis del estado del arte por ejes (E1--E6)}
  \label{tab:sintesis-ejes}
  \small
  \begin{tabular}{@{}p{0.06\textwidth}p{0.20\textwidth}p{0.22\textwidth}p{0.16\textwidth}p{0.24\textwidth}@{}}
    \toprule
    \textbf{Eje} & \textbf{Papers} & \textbf{Hallazgo principal} & \textbf{Aporte S/E/C} & \textbf{Vacío OpenAPI--BIAN} \\
    \midrule
    \textbf{E1} &
    Casas 2021; Mainas 2023; Santos y Casas 2025 &
    OpenAPI consolidado; enriquecimiento semántico posible &
    S (Mainas); contexto estándar &
    Sin vínculo con \ServiceDomain{} BIAN \\
    \textbf{E2} &
    Shvaiko 2005 &
    Taxonomía de \textit{schema matching} &
    E &
    Sin aplicación REST/banca/BIAN \\
    \textbf{E3} &
    El Bayadh 2015; Chandrasekaran 2022; Santos Sousa 2026 &
    Ontología matching + embeddings + STS &
    S; C parcial &
    Sin contratos OpenAPI ni dominio bancario \\
    \textbf{E4} &
    Schmidt 2020 &
    Dominio (DDD/MDE) $\rightarrow$ microservicios &
    Contexto dominio$\rightarrow$interfaz (no \StructuralScore{} operativo) &
    Sin métrica de alineación \\
    \textbf{E5} &
    Ehsan 2022; Golmohammadi 2024; Otieno 2026; Tanveer 2025; Vitorino 2025 &
    Testing, ML-test, seguridad &
    Delimitación (fuera de foco) &
    Orthogonal al \AlignmentScore{} \\
    \textbf{E6} &
    Mahanra Rao 2025 &
    BERT + QoS en selección de servicios &
    S parcial &
    Recomendación $\neq$ alineación BIAN \\
    \bottomrule
  \end{tabular}
\end{table}


El Anexo~B baja objetivo, técnicas, métricas y límites de cada paper. Esta sección no lo duplica.
